% Appendix G (track hw). Former Appendix D of the 438f6f16 mirror (qutrit catalog) with the qubit
% catalog of the former Sec. 3.5. The parts that restate Secs. 2.5 to 2.7 (drawing of the sequences,
% noise after each gate, scheduling) are gone. Changes of content:
% jpw-review/hw/content-changes.md, part J.
% Word lengths and gate counts: jpw-review/hw/checks/check_catalog_counts.py (.out): 216 words, 1 to 9 symbols, mean 1117/216 = 5.171; logical Clifford 16.83 one-qutrit and 3.17 two-qutrit gates; qubit: 24 words, 1 to 5 symbols, mean 61/24 = 2.542; 8.17 one-qubit and 1.50 two-qubit gates.
\section{Clifford-Group Catalogs}
\label{sec:clifford-gen}

The baseline and LRB-D draw every Clifford from a catalog of words in the primitive gates (\autoref{sub:lrb_d}, \autoref{sub:rb_baseline}). This appendix describes how the catalog of the single-qutrit Clifford group is built, how a word becomes a circuit, how the catalog was checked, and the corresponding catalog for the qubit code.

\subsection{Qutrit Catalog}
\label{app:catalog-qutrit}

Modulo global phases, the single-qutrit Clifford group $\mathcal{C}_3$ contains $216$ elements. It factors into the qutrit Pauli group and the quotient $\mathcal{C}_3 / \mathcal{P}_3$, which is isomorphic to the symplectic group $Sp(2,\mathbb{Z}_3) = SL(2,\mathbb{Z}_3)$ of order $24$~\cite{Appleby2005,Hostens2005}. The Pauli group contributes the nine operators $X^a Z^b$ with $a,b \in \mathbb{Z}_3$, and $9 \times 24 = 216$. We build the catalog in two stages, first the $24$ coset representatives and then their products with the nine Pauli operators.

\paragraph{Coset representatives.}
The quotient $\mathcal{C}_3/\mathcal{P}_3$ is generated by the images of $H$ and $S$. Two Cliffords lie in the same coset exactly when they act in the same way, by conjugation, on the eight non-identity Pauli classes $\{X, Z, XZ, X^2, Z^2, X^2Z, Z^2X, X^2Z^2\}$, so this action realizes the quotient faithfully as a permutation group on eight symbols. We traverse the Cayley graph of this group breadth first (\autoref{fig:cayley-c3modp3}). Starting from the identity, every element of the current level is multiplied on the right by $H$ and then by $S$, and an element that appears for the first time receives the word of its parent extended by the symbol of the generator. Each of the $24$ elements thus receives a word of minimal length in $H$ and $S$, and the longest of these words has five symbols. In the figure and in the text we write repeated symbols as powers, so that $SSHS$ appears as $S^2HS$, while in the circuits every symbol is a separate gate.

\begin{figure}[tb]
  \centering
  \resizebox{0.92\textwidth}{!}{%
\begin{tikzpicture}[>={Stealth[length=2.2mm,width=1.6mm]},
  cnode/.style={circle,draw=black,fill=white,inner sep=1.1pt,minimum size=7.5mm,font=\scriptsize},
  Hedge/.style={->,HCol,line width=0.75pt},
  Sedge/.style={->,SwapCol,line width=0.75pt}]
  \node[cnode] (n0) at (2.83,5.09) {$I$};
  \node[cnode] (n1) at (-1.32,6.5) {$H$};
  \node[cnode] (n2) at (2.41,0.84) {$S$};
  \node[cnode] (n3) at (1.74,6.4) {$H^2$};
  \node[cnode] (n4) at (-0.84,2.41) {$HS$};
  \node[cnode] (n5) at (-2.41,-0.84) {$SH$};
  \node[cnode] (n6) at (6.5,1.32) {$S^2$};
  \node[cnode] (n7) at (-2.49,5.26) {$H^3$};
  \node[cnode] (n8) at (1.0,2.35) {$H^2S$};
  \node[cnode] (n9) at (2.35,-1.0) {$HSH$};
  \node[cnode] (n10) at (-5.09,2.83) {$HS^2$};
  \node[cnode] (n11) at (-6.5,-1.32) {$SHS$};
  \node[cnode] (n12) at (5.09,-2.83) {$S^2H$};
  \node[cnode] (n13) at (-2.35,1.0) {$H^3S$};
  \node[cnode] (n14) at (-1.0,-2.35) {$H^2SH$};
  \node[cnode] (n15) at (5.26,2.49) {$H^2S^2$};
  \node[cnode] (n16) at (6.4,-1.74) {$HSHS$};
  \node[cnode] (n17) at (-6.4,1.74) {$SHSH$};
  \node[cnode] (n18) at (-2.83,-5.09) {$SHS^2$};
  \node[cnode] (n19) at (1.32,-6.5) {$S^2HS$};
  \node[cnode] (n20) at (0.84,-2.41) {$H^3SH$};
  \node[cnode] (n21) at (-5.26,-2.49) {$H^2SHS$};
  \node[cnode] (n22) at (2.49,-5.26) {$HSHS^2$};
  \node[cnode] (n23) at (-1.74,-6.4) {$S^2HSH$};
  \draw[Hedge] (n0) to[bend left=16] (n1);
  \draw[Sedge] (n0) to[bend left=16] (n2);
  \draw[Hedge] (n1) to[bend left=10] (n3);
  \draw[Sedge] (n1) to[bend left=16] (n4);
  \draw[Hedge] (n2) to[bend left=16] (n5);
  \draw[Sedge] (n2) to[bend left=16] (n6);
  \draw[Hedge] (n3) to[bend left=16] (n7);
  \draw[Sedge] (n3) to[bend left=16] (n8);
  \draw[Hedge] (n4) to[bend left=16] (n9);
  \draw[Sedge] (n4) to[bend left=16] (n10);
  \draw[Hedge] (n5) to[bend left=16] (n8);
  \draw[Sedge] (n5) to[bend left=16] (n11);
  \draw[Hedge] (n6) to[bend left=16] (n12);
  \draw[Sedge] (n6) to[bend left=16] (n0);
  \draw[Hedge] (n7) to[bend left=16] (n0);
  \draw[Sedge] (n7) to[bend left=16] (n13);
  \draw[Hedge] (n8) to[bend left=16] (n14);
  \draw[Sedge] (n8) to[bend left=16] (n15);
  \draw[Hedge] (n9) to[bend left=16] (n13);
  \draw[Sedge] (n9) to[bend left=16] (n16);
  \draw[Hedge] (n10) to[bend left=16] (n11);
  \draw[Sedge] (n10) to[bend left=16] (n1);
  \draw[Hedge] (n11) to[bend left=10] (n17);
  \draw[Sedge] (n11) to[bend left=16] (n18);
  \draw[Hedge] (n12) to[bend left=16] (n15);
  \draw[Sedge] (n12) to[bend left=16] (n19);
  \draw[Hedge] (n13) to[bend left=16] (n20);
  \draw[Sedge] (n13) to[bend left=16] (n17);
  \draw[Hedge] (n14) to[bend left=16] (n2);
  \draw[Sedge] (n14) to[bend left=16] (n21);
  \draw[Hedge] (n15) to[bend left=16] (n16);
  \draw[Sedge] (n15) to[bend left=16] (n3);
  \draw[Hedge] (n16) to[bend left=10] (n6);
  \draw[Sedge] (n16) to[bend left=16] (n22);
  \draw[Hedge] (n17) to[bend left=16] (n21);
  \draw[Sedge] (n17) to[bend left=16] (n7);
  \draw[Hedge] (n18) to[bend left=16] (n19);
  \draw[Sedge] (n18) to[bend left=16] (n5);
  \draw[Hedge] (n19) to[bend left=10] (n23);
  \draw[Sedge] (n19) to[bend left=16] (n20);
  \draw[Hedge] (n20) to[bend left=16] (n4);
  \draw[Sedge] (n20) to[bend left=16] (n12);
  \draw[Hedge] (n21) to[bend left=16] (n10);
  \draw[Sedge] (n21) to[bend left=16] (n23);
  \draw[Hedge] (n22) to[bend left=16] (n18);
  \draw[Sedge] (n22) to[bend left=16] (n9);
  \draw[Hedge] (n23) to[bend left=16] (n22);
  \draw[Sedge] (n23) to[bend left=16] (n14);
\end{tikzpicture}
  }
  \caption{Cayley graph of the quotient $\mathcal{C}_3/\mathcal{P}_3 \cong SL(2,\mathbb{Z}_3)$. Each of the $24$ nodes is one coset representative, labeled by the word that the breadth-first traversal assigned to it, with repeated symbols written as powers. Orange arrows denote right multiplication by $H$ and purple arrows right multiplication by $S$. The $S$ arrows close into eight directed $3$-cycles and the $H$ arrows into six directed $4$-cycles, as the orders of the two generators require, and every element is reached from the identity $I$ within five steps.}
  \label{fig:cayley-c3modp3}
\end{figure}

\paragraph{Words of the catalog.}
The catalog consists of the products $P \cdot C$, where $P = X^a Z^b$ runs over the nine Pauli operators and $C$ over the $24$ coset words. Identity factors are dropped, and the product of the two identities is the single word $I$. This gives $216$ words in the symbols $X$, $Z$, $H$, $S$ and $I$, such as $X^2Z^2S^2HSH$. A word is read as an operator product, so that its rightmost symbol acts first and its Pauli prefix last. The words contain between one and nine symbols, $5.17$ on average.

\paragraph{From a word to a circuit.}
In the baseline each symbol of a word is one primitive gate on the qutrit, and the identity element is one identity gate. Gates are never merged, so that $X^2$ is two $X$ gates and $H^3$ is three $H$ gates. In LRB-D each symbol becomes its logical circuit: $\bar{H}$ and $\bar{S}$ are the circuits of Eqs.~\eqref{eq:logical_H} and~\eqref{eq:logical_HS} (\autoref{fig: logical H}, \autoref{fig: logical S}), $\bar{X} = X_0 X_3$ and $\bar{Z} = Z_0 Z_2$ are two one-qutrit gates each, and the logical identity is one identity gate on $q_0$. Averaged over the catalog, a logical Clifford applies $16.8$ one-qutrit and $3.2$ two-qutrit gates. Because every symbol remains a separate noisy operation, the noise of a Clifford grows with the length of its word.

\paragraph{Inverses.}
For a word $C = s_1 s_2 \cdots s_\ell$ the inverse is $C^{-1} = s_\ell^{-1} \cdots s_1^{-1}$, which in time order applies $s_1^{-1}$ first and $s_\ell^{-1}$ last. The catalog stores this inverse word for each of its $216$ elements. It is written in $H^{-1}$, $S^{-1}$, $X^{-1}$ and $Z^{-1}$, has as many symbols as the forward word, and is not replaced by a shortest word. The mirror inverse looks up the stored inverse of each Clifford of the sequence and applies it like any other word. In LRB-D the inverse of a logical symbol is its logical circuit in reverse order with every gate inverted. $\bar{H}^{-1}$, for example, applies $H$ to $q_1$ and $H^{-1}$ to $q_0$, $q_2$, $q_3$ and $q_4$ and then $\mathrm{SWAP}_{2,3}$.

\paragraph{Checks.}
The catalog and all circuits are built as \texttt{sdim} circuits, with the gates of Eq.~\eqref{eq:conventions}. We checked the catalog against these conventions by converting the single-qutrit circuit of every word to its unitary matrix. The $216$ words give $216$ distinct unitaries modulo a global phase, which coincide with the closure of $\{H, S\}$ under multiplication, and every word followed by its stored inverse gives the identity up to a phase. The same unitaries reproduce the $24$ coset words of the traversal and the $48$ edges of \autoref{fig:cayley-c3modp3}.

\subsection{Qubit Catalog}
\label{app:catalog-qubit}

% from "method.tex", line 16 of the 438f6f16 mirror
% Qubit code and catalog: FoldedQubitDetectionCode and QubitCliffordLibrary (worktree qubit-lrbd-sdim140, src/lrb/code_definitions.py), FoldedQubitConventionalSitesDetectionCode (worktree conventional-sites); checks in tests/test_folded_qubit_code.py and tests/test_conventional_sites_code.py of those worktrees.
The catalog of the $24$-element single-qubit Clifford group consists of the products of the Pauli operators $I$, $X$, $Z$ and $XZ$ with the six words $I$, $H$, $S$, $HS$, $SH$ and $HSH$, with identity factors dropped as for the qutrit. Its words contain between one and five symbols, $61/24 = 2.54$ on average, and a logical Clifford of the qubit code applies $8.2$ one-qubit and $1.5$ two-qubit gates on average. Inverses and circuits follow the rules of the qutrit catalog. The encoded circuits of all $24$ words and of their inverses act on the code space as the intended single-qubit Cliffords, and in a noiseless simulation of either layout every Pauli error of weight one on a data qubit gives a nontrivial syndrome in a check round.
